% !PRINTED 349
\weberpage{349}
%<<CONT
gerade sein muss, gleich 6. Zwei reelle oder zwei conjugirt imaginäre Wurzeln
$x_1, x_2$ gehören immer einem Tripel an, dessen dritte Wurzel reell ist; dies
ergiebt sich aus der Gleichung:

\[
x_3 = \Theta (x_1, x_2) = \Theta (x_2, x_1).
\]

Wir wollen demnach ein Tripel, was zwei reelle oder zwei conjugirt imaginäre Wurzeln
enthält, ein \gesperrt{reelles Tripel} nennen. Wenn in einem Tripel irgend imaginäre
Wurzeln vorkommen, so bilden auch die conjugirt imaginären Wurzeln ein Tripel, und
zwei solche Tripel sollen \gesperrt{conjugirt imaginäre Tripel} heissen.

Es giebt dann drei Möglichkeiten:

1. \gesperrt{Lauter reelle Wurzeln.}

2. \gesperrt{Eine reelle und acht imaginäre Wurzeln.}

In diesem Falle müssen die vier Paare conjugirt imaginärer Wurzeln vier Tripel
bestimmen, die alle dieselbe reelle Wurzel als dritte enthalten. Bezeichnen wir
diese reelle Wurzel mit $(0,0)$, so müssen

\[
(1,2)\ (2,1),\quad (1,1)\ (2,2),\quad (1,0)\ (2,0),\quad (0,1)\ (0,2)
\]

die vier Paare conjugirter Wurzeln sein.

3. \gesperrt{Drei reelle und sechs imaginäre Wurzeln.}

Die drei reellen Wurzeln müssen hier ein Tripel bilden. Nehmen wir für die reellen
Wurzeln

\weq{1}{(0,0)\ (1,1)\ (2,2),}

so können in den beiden Reihen

\weq{2}{\begin{aligned} &(1,2)\ (0,2)\ (1,0) \\ &(2,1)\ (2,0)\ (0,1) \end{aligned}}

conjugirt imaginäre Wurzeln nicht in derselben Reihe vorkommen, weil sonst die
dritte Wurzel der betreffenden Reihe reell wäre.

Es ist zu zeigen, dass die durch die conjugirten Paare bestimmten drei reellen
Tripel zusammen alle drei reellen Punkte enthalten müssen.

Nehmen wir nämlich an, unter den drei Tripeln

\[
(0,0)\ (1,2)\ (2,1),\quad (0,0)\ (2,0)\ (1,0),\quad (0,0)\ (0,1)\ (0,2)
\]

kommen zwei reelle vor, so muss auch das dritte reell sein, d.\,h. es müssen

\[
(1,2)\ (2,1),\quad (2,0)\ (1,0),\quad (0,1)\ (0,2)
\]
%CONT>>
